% Generated draft: requires mathematical and source review.
\documentclass[11pt]{article}
\usepackage[margin=25mm]{geometry}
\usepackage{fontspec}
\setmainfont{Latin Modern Roman}
\usepackage{amsmath,amssymb,mathtools,unicode-math}
\setmathfont{Latin Modern Math}
\usepackage{enumitem,needspace,xurl,hyperref}
\setlength{\parindent}{0pt}
\setlength{\parskip}{5pt}
\setlength{\emergencystretch}{3em}
\newcommand{\N}{\mathbb{N}}
\newcommand{\Z}{\mathbb{Z}}
\newcommand{\Q}{\mathbb{Q}}
\newcommand{\R}{\mathbb{R}}
\newcommand{\C}{\mathbb{C}}
\newcommand{\F}{\mathbb{F}}
\newcommand{\A}{\mathbb{A}}
\newcommand{\PP}{\mathbb{P}}
\newcommand{\E}{\mathbb{E}}
\newcommand{\eps}{\varepsilon}
\newcommand{\dd}{\,\mathrm{d}}
\begin{document}
\section*{UnsolvedMath ID: 20002685}
\textbf{Source problem code:} catalogue entry
\subsection{Definitions and mathematical statement}
Question 25 asks whether two local-growth spaces $K(T,r)$ and $E(T,r)$ agree for $T$ in a $\mathrm{II}_1$ factor, suggesting a DT quasinilpotent operator as a possible counterexample. The supplied aid says the rendered source itself has malformed definitions: $r$ is unused, a limsup is followed by an extra arrow to zero, and the second set contains an unquantified approximating sequence. Thus there is no fully coherent literal statement.

The aid explicitly reconstructs the following Haagerup--Schultz-type comparison, not a verified verbatim recovery. Let $M$ act faithfully and normally on a Hilbert space $H$ (for example $L^2(M,\operatorname{tr})$), $T\in M$, and $r\geq0$. Define
\[
K_{\mathrm{app}}(T,r)=\{\xi\in H:\exists\xi_n\to\xi\text{ in norm},\ \limsup_{n\to\infty}\|T^n\xi_n\|^{1/n}\leq r\},
\]
\[
E_{\mathrm{pt}}(T,r)=\{\xi\in H:\limsup_{n\to\infty}\|T^n\xi\|^{1/n}\leq r\}.
\]
The approximation is exponent dependent in the first definition; the vector in the second is fixed. Always $E_{\mathrm{pt}}\subseteq K_{\mathrm{app}}$. The aid reports a negative answer to universal equality under this reconstruction.

Its counterexample exists inside every $\mathrm{II}_1$ factor. Choose nonzero mutually orthogonal projections $p_m$ summing strongly to $I$, and in each corner $p_mMp_m$ choose an $m$-step nilpotent shift $S_m$ of norm at most one from a unital matrix-unit system. Put $T=\sum_mS_m$ in the strong sense. The aid asserts
\[
\|T\|=1,\qquad\sigma(T)=\{z:|z|\leq1\},\qquad
\mu_T=\delta_0,\qquad |T^n|^{1/n}\to0\text{ strongly},
\]
where $\mu_T$ denotes Brown measure. For every vector, its finite-block approximants $\xi_n=\sum_{m\leq n}p_m\xi$ satisfy $T^n\xi_n=0$, so $K_{\mathrm{app}}(T,r)=H$ for all $r\geq0$. Choose orthogonal unit vectors $\eta_m\in p_mH$ with $\|S_m^j\eta_m\|=1$ for $0\leq j<m$, and let $\xi=\sum_m m^{-1}\eta_m$. Then
\[
\|T^n\xi\|^2=\sum_{m>n}m^{-2},\qquad
\lim_{n\to\infty}\|T^n\xi\|^{1/n}=1.
\]
Consequently $E_{\mathrm{pt}}(T,r)$ is proper for every $0\leq r<1$, whereas both spaces equal $H$ for $r\geq1$. Brown measure concentrated at zero and strong convergence of the root moduli must not be confused with spectral quasinilpotence, since this example has spectral radius one.

The aid also reports the exact criterion
\[
K_{\mathrm{app}}(T,r)=E_{\mathrm{pt}}(T,r)
\quad\Longleftrightarrow\quad
r\bigl(T|_{K_{\mathrm{app}}(T,r)}\bigr)\leq r,
\]
using that $K_{\mathrm{app}}$ is a closed invariant subspace, and with the zero-subspace case interpreted trivially. The radius on the left of the inequality is operator spectral radius; the scalar on the right is the growth threshold. In particular a genuinely spectrally quasinilpotent operator has $E_{\mathrm{pt}}(T,0)=H$, so the original DT suggestion requires checking which meaning of quasinilpotence and which definitions were intended. It is not needed for the displayed block-shift counterexample.
\subsection{Short English statement}
The source's two definitions are malformed. Under the aid's explicit reconstruction, exponent-dependent approximants can remove growth that a fixed vector retains. A block sum of nilpotent shifts gives full approximating space but a proper fixed-vector space below threshold one, despite Brown measure concentrated at zero.
\subsection{Sources}
\begin{itemize}
\item Local UnsolvedMath registry, permanent ID $20002685$, \texttt{catalogue entry}; \texttt{aim-probability-notes.json}, \texttt{notes[126]}.
\item American Institute of Mathematics, \emph{Free analysis}, Question 25, \url{https://aimath.org/WWN/freeanalysis/freeanalysis.pdf}. The source suggests a DT quasinilpotent operator, and the aid reports that the malformed definitions occur in the rendered PDF itself.
\item Haagerup--Schultz is the supplied attribution for the coherent local-growth-space reconstruction; no title, theorem number, or authoritative full citation is supplied. The block-shift construction and spectral-radius criterion are supplied through the research aid.
\item Research attempt 1: counterexample, problem refuted at run, rigor label ``proved,'' novelty confidence low. Registry classification solved in literature, status solved, statement status reconstructed/unverified. The refutation concerns the explicitly reconstructed universal equality.
\end{itemize}
\end{document}
